\begin{landscape}
    \begin{longtable}{p{2cm}p{2cm}p{2cm}p{6cm}p{2cm}p{2cm}}
        \caption{心理学研究92巻で用いられた尺度項目など}
        \label{JPA92Scales}                                                                                                                                      \\
        \hline
        出典                       & 尺度名など                            & 項目数 & 項目例                                                 & 因子名                  & $\alpha$係数 \\ \hline
        \hline \endhead
        \citet{Suzuki202192}     & 動機づけ尺度                           & 52  & できないことができるようになるとうれしいから                              & 内的調整                 & 0.820   \\
                                 &                                  & 52  & 集団活動や礼儀について学べるから                                    & 同一化調整                & 0.800   \\
                                 &                                  & 52  & サボっていると思われたくないから                                    & 取り入れ調整               & 0.810   \\
                                 &                                  & 52  & 部活に参加しないと罰を受けるから                                    & 外的調整                 & 0.810   \\
                                 &                                  & 52  & 活動に参加したくない                                          & 無調整                  & 0.840   \\
                                 & 部活動への傾倒尺度                        & 4   & 部活には自主的に参加している                                      & 傾倒                   & 0.840   \\
                                 & 基本的心理欲求充足尺度                      & 12  & 私は，能力のある人間だと感じている                                   & 有能感                  & 0.730   \\
                                 &                                  & 12  & 私は，周りの人と有効な関係を築いていると思う                              & 関係性                  & 0.780   \\
                                 &                                  & 12  & 私は自分の意見や考えを自分から自由に言えていると思う                          & 自律性                  & 0.740   \\
                                 & 動機づけ尺度                           & 52  & できないことができるようになるとうれしいから                              & 内的調整                 & 0.870   \\
                                 &                                  & 52  & 集団活動や礼儀について学べるから                                    & 同一化調整                & 0.840   \\
                                 &                                  & 52  & サボっていると思われたくないから                                    & 取り入れ調整               & 0.820   \\
                                 &                                  & 52  & 部活に参加しないと罰を受けるから                                    & 外的調整                 & 0.760   \\
                                 &                                  & 52  & 活動に参加したくない                                          & 無調整                  & 0.840   \\
                                 & リーダーシップ尺度                        & 17  & 気まずい雰囲気があると解きほぐす                                    & 関係調整・統率              & 0.920   \\
                                 &                                  & 17  & 活動内容に関する知識が豊富である                                    & 技術指導                 & 0.850   \\
                                 &                                  & 17  & 厳しく命令したり注意したりする                                     & 規範的指導                & 0.840   \\
        \citet{Kambara202192}    & バイアスについての説明文                     & 7   & 自分の信念にあう情報は受け入れるが，自分の信念に矛盾する情報は否定する。                &                      & 0.710   \\
                                 & 驚いた程度                            & 1   & カテゴリ直接選択                                            &                      &         \\
        \citet{Kawasaki202192}   & 病理的自己愛目録日本語版                     & 52  & 私が人の手伝いをするのは，自分が良い人間であることをわかってもらうためだ                & 誇大性                  &         \\
                                 &                                  & 52  & 私はよく，自分の偉業が認められるという空想をする                            & 誇大妄想                 &         \\
                                 &                                  & 52  & 誰に対しても，私の思うように物事を信じさせることができる                        & 搾取性                  &         \\
                                 &                                  & 52  & 私が言ったこと・したことに，興味を持ってもらえないとイライラする                    & 脆弱性                  &         \\
                                 &                                  & 52  & 心の中で感じている自分の弱さを他人に見せるのは耐えがたいことだ                     & 自己隠蔽                 &         \\
                                 &                                  & 52  & 人が自分の存在に気がついてくれないと，がっかりしてしまう                        & 随伴的自尊感情              &         \\
                                 &                                  & 52  & 私がしてあげたことが感謝されないのではないかと思って，人を避けてしまうことがある            & 脱価値化                 &         \\
        \citet{Namatame202192}   & 日本語版幸せへのおそれ尺度                    & 5   & 喜びの後にはたいてい悲しみがやってくるので，私はあまり喜びすぎないようにしている            & 幸せへのおそれ              & 0.880   \\
                                 & 日本語版幸せの壊れやすさ                     & 4   & どんな時でも何かが起こってしまって，幸せをあっけなく失うかもしれない                  & 幸せの壊れやすさ             & 0.900   \\
        \citet{Ono202192}        & 感情反応を測定する項目                      & 17  & やる気がわく                                              & ポジティブ感情              & 0.920   \\
                                 &                                  & 17  & 警戒する                                                & 不信感                  & 0.890   \\
                                 &                                  & 17  & 怒りを感じる                                              & 苛立ち                  & 0.770   \\
                                 & 魅力を測定する項目                        & 1   & メッセージカードの送り手は魅力的だと思う                                & 魅力評定                 &         \\
                                 & 感情反応を測定する項目                      & 17  & やる気がわく                                              & ポジティブ感情              & 0.810   \\
                                 &                                  & 17  & 警戒する                                                & 不信感                  & 0.790   \\
                                 &                                  & 17  & 怒りを感じる                                              & 苛立ち                  & 0.870   \\
        \citet{Niikuni202192}    & 主体感を測定する尺度                       & 17  &                                                     & 精神的活動における主体の誤帰属      &         \\
                                 &                                  & 17  & 自分の身体を思うように動かせないと感じることがある                           & 身体的活動における自己身体の制御不能性  &         \\
                                 &                                  & 17  & 周りに協調するよりも自分の主張を通すことがある                             & 社会的活動における自己の主張性      &         \\
        \citet{Toyoda202192}     & 文章刺激に対する評価                       & 1   & 先ほどの場面についての想像は，どのぐらい鮮明でしたか                          & 想像の鮮明さ               &         \\
                                 &                                  & 1   & ◯さんはどのような気持ちだと思いますか                                 & 相手の感情                &         \\
                                 &                                  & 1   & ◯さんについての文章を読んで，あなたはどのような気持ちになりましたか                  & 自分の感情                &         \\
                                 &                                  & 1   & ◯さんに対して，あなたはどのぐらい助けようと思いましたか                        & 援助意図                 &         \\
                                 &                                  & 1   & あなたがもし，◯さんと同じ状況だとしたら，他者に助けを求めますか                    & 援助希求度                &         \\
                                 &                                  & 1   & あなたは過去に，◯さんと同じような状況の人物に遭遇したことがありますか                 & 遭遇頻度                 &         \\
                                 &                                  & 1   & あなたは過去に，◯さんと同じような状況になったことがありますか                     & 経験頻度                 &         \\
                                 &                                  & 1   & 先ほどの人物はどのような気持ちだと思いますか                              & 相手の感情                &         \\
                                 &                                  & 1   & 先ほどの人物についての文章を読んで，あなたはどのような気持ちになりましたか               & 自分の感情                &         \\
                                 &                                  & 1   & 先ほどの人物に対して，あなたはどのぐらい助けようと思いましたか                     & 援助意図                 &         \\
                                 &                                  & 1   & あなたがもし，先ほどの人物と同じ状況だとしたら，他者に助けを求めますか                 & 援助希求度                &         \\
                                 &                                  & 1   & あなたは過去に，先ほどの人物と同じような状況の人物に遭遇したことがありますか              & 遭遇頻度                 &         \\
                                 &                                  & 1   & あなたは過去に，先ほどの人物と同じような状況になったことがありますか                  & 経験頻度                 &         \\
        \citet{Hirose202192}     & 日本語版SIS                          & 32  & 私は悲しくなる                                             & 情緒覚醒                 & 0.770   \\
                                 &                                  & 32  & 私の一日は台無しになる                                         & 情緒統制不全               & 0.810   \\
                                 &                                  & 32  & 私は自分の気持ちを隠そうとする                                     & 葛藤からの回避              & 0.820   \\
                                 &                                  & 32  & 私はどちらか，または両方に対してかわいそうに感じる                           & 葛藤への関与               & 0.750   \\
                                 &                                  & 32  & 家族はまだうまくやっていける                                      & 建設的な家族表象             & 0.860   \\
                                 &                                  & 32  & 私は家族の将来が心配になる                                       & 破壊的な家族表象             & 0.790   \\
                                 &                                  & 32  & 私は板挟みになっているように感じる                                   & 巻き込まれ表象              & 0.800   \\
                                 & 破壊的IC                            & 20  & 両親はよく喧嘩をする                                          & 葛藤の激しさ               & 0.810   \\
                                 &                                  & 20  & 両親はけんかをしても，すぐに仲直りする                                 & 葛藤の持続性               & 0.800   \\
                                 &                                  & 20  & 両親はけんかをしても，大きな声を出さずに話し合う                            & 葛藤の解決                & 0.740   \\
                                 & ICのおそれ・身体反応                      & 5   & 両親がけんかをしていると，不安になる                                  & ICのおそれ・身体反応          & 0.880   \\
                                 & 適応問題                             & 33  &                                                     & 不安・抑うつ               & 0.850   \\
                                 &                                  & 33  &                                                     & 攻撃的行動                & 0.860   \\
        \citet{Amai202192}       & 中学生を対象とした情緒的援助期待尺度               & 64  & とにかくただ話を聞いてほしい                                      & 受容期待                 & 0.760   \\
                                 &                                  & 64  & 今の自分が客観的にどう見えるか知りたい                                 & 再解釈期待                & 0.610   \\
                                 &                                  & 64  & 自分はいい人間だと思わせてほしい                                    & 正当化期待                & 0.740   \\
                                 &                                  & 64  & がんばれと言ってほしい                                         & 楽観的期待                & 0.730   \\
                                 &                                  & 64  & 一緒に何か夢中になることをして辛い状況を忘れたい                            & 気晴らし期待               & 0.780   \\
                                 & 悩みの認知評価                          & 14  & 解決するための方法がわかっている                                    & 統制可能性                &         \\
                                 &                                  & 14  & 自分を傷つけることだと思う                                       & 影響性                  &         \\
                                 & 援助要請行動                           & 8   & 誰かから同情や理解を得る                                        & 情緒的サポートの利用           &         \\
                                 &                                  & 8   & 何をすべきか誰かからアドバイスを得ようとする                              & 道具的サポートの利用           &         \\
                                 & 中学生を対象とした情緒的援助期待尺度               & 23  & きちんと話を聞いてほしい                                        & 受容期待                 & 0.730   \\
                                 &                                  & 23  & 別な見方を教えてほしい                                         & 再解釈期待                & 0.750   \\
                                 &                                  & 23  & 間違ってないよと言ってほしい                                      & 正当化期待                & 0.820   \\
                                 &                                  & 23  & そんなこと気にするなと笑い飛ばしてほしい                                & 楽観的期待                & 0.640   \\
                                 &                                  & 23  & 全く別の話をして話題から意識をそらしたい                                & 気晴らし期待               & 0.750   \\
                                 & 援助要請行動                           & 8   & 誰かから同情や理解を得る                                        & 情緒的サポートの利用           &         \\
                                 &                                  & 8   & 何をすべきか誰かからアドバイスを得ようとする                              & 道具的サポートの利用           &         \\
                                 & 中学生を対象とした情緒的援助期待尺度               & 23  & きちんと話を聞いてほしい                                        & 受容期待                 & 0.720   \\
                                 &                                  & 23  & 別な見方を教えてほしい                                         & 再解釈期待                & 0.830   \\
                                 &                                  & 23  & 間違ってないよと言ってほしい                                      & 正当化期待                & 0.840   \\
                                 &                                  & 23  & そんなこと気にするなと笑い飛ばしてほしい                                & 楽観的期待                & 0.720   \\
                                 &                                  & 23  & 全く別の話をして話題から意識をそらしたい                                & 気晴らし期待               & 0.790   \\
                                 & 援助評価                             & 23  & どうすればいいかがはっきりした                                     & 問題状況の改善              &         \\
                                 &                                  & 23  & どうすればいいか余計にわからなくなった                                 & 退所の混乱                &         \\
                                 &                                  & 23  & 自分は一人じゃないんだと思った                                     & 他者からの支えの知覚           &         \\
                                 &                                  & 23  & 自分が他の人に頼りすぎていると思った                                  & 他者への依存               &         \\
                                 & 愛着                               & 30  & 人に対して自分のイメージを悪くしないか恐れている                            & 見捨てられ不安              &         \\
                                 &                                  & 30  & 他の人との間に壁を作っている                                      & 親密性の回避               &         \\
        \citet{Kanemasa202192}   & 愛着スタイル尺度                         & 15  &                                                     & 愛着不安                 & 0.930   \\
                                 &                                  & 15  &                                                     & 愛着不安                 & 0.940   \\
                                 & パートナーへの間接的暴力加害                   & 6   & 相手の行動を制限することが                                       & DaV加害                & 0.900   \\
                                 &                                  & 6   & 相手の行動を制限することが                                       & DaV加害                & 0.910   \\
                                 &                                  & 6   & 相手の行動を制限することが                                       & DaV被害                & 0.900   \\
                                 & 一般的他者版ECR                        & 18  &                                                     & 愛着不安                 & 0.940   \\
                                 & パートナーへの間接的暴力加害                   & 6   & 相手の行動を制限することが                                       & DaV加害                & 0.940   \\
                                 &                                  & 6   & 相手の行動を制限することが                                       & DaV加害                & 0.900   \\
                                 &                                  & 6   & 相手の行動を制限することが                                       & DaV被害                & 0.900   \\
        \citet{Moroi202192}      & 詐欺電話に関する測定尺度                     & 6   & 詐欺電話がかかってきたことを誰かに相談すると思いますか                         & 相談行動意図               &         \\
                                 &                                  & 3   & （あなたが詐欺電話について相談した場合，配偶者は）電話を詐欺だと見抜くことができると思いますか」    & 相談相手としての信頼(能力の予期)・妻  &         \\
                                 &                                  & 3   & （あなたが詐欺電話について相談した場合，配偶者は）電話を詐欺だと見抜くことができると思いますか」    & 相談相手としての信頼(能力の予期)・夫  & 0.780   \\
                                 &                                  & 3   & （あなたが詐欺電話について相談した場合，配偶者は）誠実に対応すると思いますか」             & 相談相手としての信頼(誠実さの予期)・妻 &         \\
                                 &                                  & 3   & （あなたが詐欺電話について相談した場合，配偶者は）誠実に対応すると思いますか」             & 相談相手としての信頼(誠実さの予期)・夫 &         \\
                                 &                                  & 3   & （あなたが詐欺電話について相談した場合，配偶者は）あなたの気持ちを分かってくれると思いますか」     & 相談相手としての信頼(共感の予期)・妻  &         \\
                                 &                                  & 3   & （あなたが詐欺電話について相談した場合，配偶者は）あなたの気持ちを分かってくれると思いますか」     & 相談相手としての信頼(共感の予期)・夫  &         \\
                                 &                                  & 4   & あなたは詐欺電話がかかって来た時に，自分がお金を騙し取られる被害にあう可能性がどの程度あると思いますか & 被害リスク認知・妻            &         \\
                                 &                                  & 4   & あなたは詐欺電話がかかって来た時に，自分がお金を騙し取られる被害にあう可能性がどの程度あると思いますか & 被害リスク認知・夫            &         \\
                                 & 夫婦関係に関する測定尺度                     & 19  & うれしかったこと，たのしかったことについて                               & コミュニケーション・妻          & 0.910   \\
                                 &                                  & 19  & うれしかったこと，たのしかったことについて                               & コミュニケーション・夫          & 0.910   \\
                                 &                                  & 19  & あなたは，配偶者との関係についてどの程度満足していますか                        & 関係満足度・妻              &         \\
                                 &                                  & 19  & あなたは，配偶者との関係についてどの程度満足していますか                        & 関係満足度・夫              &         \\
                                 &                                  & 19  & 配偶者と話し合おうとする                                        & 問題対処法略の夫婦関係内アプローチ・妻  & 0.860   \\
                                 &                                  & 19  & 配偶者と話し合おうとする                                        & 問題対処法略の夫婦関係内アプローチ・夫  & 0.850   \\
                                 &                                  & 19  & 配偶者とは別に普段より社会的な催しに参加する                              & 問題対処法略の夫婦関係外アプローチ・妻  & 0.710   \\
                                 &                                  & 19  & 配偶者とは別に普段より社会的な催しに参加する                              & 問題対処法略の夫婦関係外アプローチ・夫  & 0.760   \\
        \citet{Mizuno202192}     & セルフコンパッション                       & 26  &                                                     & セルフ・コンパッション          & 0.760   \\
                                 & 新職業性ストレス簡易調査票短縮版                 & 6   &                                                     & 上司からのサポート            & 0.750   \\
                                 &                                  & 6   &                                                     & 同僚からのサポート            & 0.830   \\
                                 &                                  & 17  &                                                     & 職務ストレッサー             & 0.630   \\
                                 & バーンアウト傾向                         & 17  &                                                     & 情緒的消耗感               & 0.820   \\
                                 &                                  & 17  &                                                     & 脱人格化                 & 0.800   \\
                                 &                                  & 17  &                                                     & 個人的達成感の後退            & 0.810   \\
        \citet{Tabata202192}     & インターネット上での情報プライバシー               & 20  & 最近1ヶ月のうちに，以下の情報を含めたツイート，リツイートをしたことがどのくらいありましたか      & 情報公開回数               & 0.910   \\
                                 & 他者懸念                             & 7   & 他人にプライベートな質問をされたくない                                 & 自己意識・維持行動            & 0.760   \\
                                 &                                  & 4   &                                                     & 他者意識                 & 0.690   \\
                                 &                                  & 4   &                                                     & 他者維持行動               & 0.740   \\
                                 &                                  & 5   & 総合的に考えて，他人についての情報を他の利用者と共有することは，深刻なプライバシーの問題を引き起こす  & Twitter上での他者懸念       & 0.890   \\
        \citet{Ishikawa202192}   & 気分状態                             & 24  & いらいらしている                                            & 緊張と興奮                &         \\
                                 &                                  & 24  & 気持ちが滅入っている                                          & 抑うつ感                 &         \\
                                 &                                  & 24  & 何か物足りない                                             & 不安感                  &         \\
        \citet{Ikegmai202192}    & 音楽聴取の心理的機能                       & 133 & 音楽は自分自身の道を見つけてくれるのを助けてくれるから                         & 自己認識                 &         \\
                                 &                                  & 133 & 音楽は私の機敏を明るくすることができるから                               & 感情調節                 &         \\
                                 &                                  & 133 & 音楽は友人たちとの話題にできるものだから                                & コミュニケーション            &         \\
                                 &                                  & 133 & 私は周りの雑音を音楽でかき消したいから                                 & 道具的活用                &         \\
                                 &                                  & 133 & 大音量で聴くのが好きだから                                       & 身体性                  &         \\
                                 &                                  & 133 & 音楽は私の信仰心を支えてくれるから                                   & 社会的距離調節              &         \\
                                 &                                  & 133 & 音楽は私の悲しみに寄り添ってくれるから                                 & 慰め                   &         \\
        \citet{Akazawa202192}    & アサーション                           & 16  & 好きな人には素直に愛情や好意を示す                                   & 関係形成                 & 0.785   \\
                                 &                                  & 16  & 買った商品に欠陥があったら交換してもらう                                & 説得交渉                 & 0.757   \\
                                 &                                  & 16  & 好きな人には素直に愛情や好意を示す                                   & 関係形成                 & 0.767   \\
                                 &                                  & 16  & 買った商品に欠陥があったら交換してもらう                                & 説得交渉                 & 0.786   \\
                                 &                                  & 16  & 好きな人には素直に愛情や好意を示す                                   & 関係形成                 & 0.642   \\
                                 &                                  & 16  & 買った商品に欠陥があったら交換してもらう                                & 説得交渉                 & 0.610   \\
                                 & 視点取得                             & 5   &                                                     & 視点取得                 & 0.912   \\
                                 &                                  & 5   &                                                     & 視点取得                 & 0.944   \\
                                 &                                  & 5   &                                                     & 視点取得                 & 0.929   \\
                                 & 暴力観                              & 5   & けがをしない強さで叩く                                         & 暴力観                  & 0.809   \\
                                 &                                  & 5   & けがをしない強さで叩く                                         & 暴力観                  & 0.890   \\
                                 &                                  & 5   & けがをしない強さで叩く                                         & 暴力観                  & 0.771   \\
        \citet{Toyama202192}     & 制御焦点                             & 16  & どうやったら自分の目標や希望を叶えられるか，よく想像することがある                   & 利得接近志向尺度             & 0.730   \\
                                 &                                  & 16  & どうやったら失敗を防げるかについて，よく考える                             & 損失回避志向尺度             & 0.810   \\
                                 & エンゲージメント                         & 14  & 課題に一生懸命に取り組んだ                                       & 行動的エンゲージメント          & 0.890   \\
                                 &                                  & 14  & 課題は楽しかった                                            & 感情的エンゲージメント          & 0.930   \\
                                 &                                  & 14  & 課題に没頭していた                                           & 状態的エンゲージメント          & 0.800   \\
        \citet{Togo202192}       & 子の運動に対する養育態度尺度                   & 10  & 子どもと一緒に運動する                                         & 共同活動                 & 0.810   \\
                                 &                                  & 10  & 試合(あるいは進級試験等)について反省させる                              & 支配的対応                & 0.780   \\
                                 &                                  & 10  & 運動したり試合をしたりすることは，子どもにとって良いことだよとつたえる                 & 運動の価値伝授              & 0.830   \\
                                 & 子供への対応                           & 17  & 次は，どんな目的をもって練習する?                                   & 自律的対応                & 0.870   \\
                                 &                                  & 17  & これだから負けるのよ                                          & 支配・過干渉               & 0.810   \\
                                 &                                  & 17  & 今日もよく頑張ったね                                          & 肯定的評価                & 0.720   \\
                                 &                                  & 17  & 昨日のテレビの話など試合に関係ない話題をする                              & 間接的支援                & 0.820   \\
                                 &                                  & 17  & 次は，どんな目的をもって練習する?                                   & 自律的対応                & 0.880   \\
                                 &                                  & 17  & これだから負けるのよ                                          & 支配・過干渉               & 0.850   \\
                                 &                                  & 17  & 今日もよく頑張ったね                                          & 肯定的評価                & 0.690   \\
                                 &                                  & 17  & 昨日のテレビの話など試合に関係ない話題をする                              & 間接的支援                & 0.820   \\
        \citet{Sakakibara202192} & マスク着用の理由                         &     & あなたがもし新型コロナウイルスに感染したら，症状は深刻なものになると思いますか             &                      &         \\
                                 &                                  &     & マスクは，自分が周りからウイルスを移されないために着用している                     &                      &         \\
        \citet{Miyazaki202192}   & 他者に見られる不安                        & 20  & 人前で文字を書かなければならない時，不安になる                             &                      & 0.940   \\
                                 & 対人交流不安                           & 20  & あまり知らない人に会った時，あいさつするかどうか迷う                          &                      & 0.940   \\
                                 & 他者に見られる不安                        & 20  & 人前で文字を書かなければならない時，不安になる                             &                      & 0.950   \\
                                 & 対人交流不安                           & 20  & あまり知らない人に会った時，あいさつするかどうか迷う                          &                      & 0.940   \\
                                 & 他者に見られる不安                        & 20  & 人前で文字を書かなければならない時，不安になる                             &                      & 0.930   \\
                                 & 対人交流不安                           & 20  & あまり知らない人に会った時，あいさつするかどうか迷う                          &                      & 0.940   \\
                                 & 他者に見られる不安                        & 20  & 人前で文字を書かなければならない時，不安になる                             &                      & 0.930   \\
                                 & 対人交流不安                           & 20  & あまり知らない人に会った時，あいさつするかどうか迷う                          &                      & 0.950   \\
                                 & 特性不安                             & 20  & たのしい                                                &                      & 0.930   \\
                                 & 感染脆弱意識尺度                         & 15  & 風邪やインフルエンザなどにとても感染しやすい                              & 易感染性                 & 0.810   \\
                                 &                                  & 15  & 誰かと握手した後は手を洗いたくなる                                   & 感染嫌悪                 & 0.750   \\
                                 & 特性不安                             & 20  & たのしい                                                &                      & 0.920   \\
                                 & 感染脆弱意識尺度                         & 15  & 風邪やインフルエンザなどにとても感染しやすい                              & 易感染性                 & 0.840   \\
                                 &                                  & 15  & 誰かと握手した後は手を洗いたくなる                                   & 感染嫌悪                 & 0.740   \\
                                 & 特性不安                             & 20  & たのしい                                                &                      & 0.920   \\
                                 & 感染脆弱意識尺度                         & 15  & 風邪やインフルエンザなどにとても感染しやすい                              & 易感染性                 & 0.790   \\
                                 &                                  & 15  & 誰かと握手した後は手を洗いたくなる                                   & 感染嫌悪                 & 0.720   \\
                                 & 特性不安                             & 20  & たのしい                                                &                      & 0.920   \\
                                 & 感染脆弱意識尺度                         & 15  & 風邪やインフルエンザなどにとても感染しやすい                              & 易感染性                 & 0.800   \\
                                 &                                  & 15  & 誰かと握手した後は手を洗いたくなる                                   & 感染嫌悪                 & 0.690   \\
                                 & 特性不安                             & 20  & たのしい                                                &                      & 0.910   \\
                                 &                                  & 15  & 風邪やインフルエンザなどにとても感染しやすい                              & 易感染性                 & 0.820   \\
                                 & 感染脆弱意識尺度                         & 15  & 誰かと握手した後は手を洗いたくなる                                   & 感染嫌悪                 & 0.740   \\
        \citet{Kamatani202192}   & 黒色マスク着用者に対する信念                   & 1   & 男性の顔の魅力は，黒い衛生マスクを着用することで高まると思いますか                   & 魅力                   &         \\
                                 &                                  & 1   & 黒い衛生マスクを着用する人物についてどう感じますか                           & 健康さ                  &         \\
                                 &                                  & 1   & 黒い衛生マスクを着用する人物についてどう感じますか                           & ファッション性              &         \\
                                 &                                  & 1   & 黒い衛生マスクを着用する人物についてどう感じますか                           & イメージの良さ              &         \\
        \citet{Yamamoto202192}   & 感染者の特性に対するステレオタイプ的認知             & 20  & 責任感のあるー責任感のない                                       & 社会性                  & 0.902   \\
                                 &                                  & 20  & 無気力なー意欲的な                                           & 活動性                  & 0.755   \\
                                 & 感染予防行動に対する認知                     & 10  & 外出の際はマスクを着用しない                                      & 感染予防行動の欠如            & 0.972   \\
                                 & 感染脆弱意識尺度                         & 15  & 他の人よりも，病気にかかりやすい方だ                                  & 易感染性                 & 0.737   \\
                                 &                                  & 15  & 前に使った人から何かうつりそうなので，公衆電話は使いたくない                      & 感染嫌悪                 & 0.610   \\
                                 & 感染予防行動に対する認知                     & 10  & 外出の際はマスクを着用する                                       & 感染予防行動               & 0.760   \\
        \citet{Iida202192}       & 経済的負担感                           & 1   & 新型コロナウイルスが発生する前の，あなた自身の経済状態について，あなたはどのように感じていますか    &                      &         \\
                                 & 孤独感                              & 3   &                                                     &                      &         \\
                                 & 精神的健康                            & 6   &                                                     & 抑うつ                  &         \\
                                 &                                  & 7   &                                                     & 不安                   &         \\
        \citet{Nagai202192}      & 主体的な学修態度尺度                       & 9   &                                                     &                      &         \\
                                 &                                  & 9   &                                                     &                      &         \\
        \citet{Nakao202192}      & アタッチメント                          & 9   &                                                     & 不安                   & 0.780   \\
                                 &                                  & 9   &                                                     & 回避                   & 0.770   \\
                                 & 孤独感                              & 20  &                                                     & 孤独感                  & 0.890   \\
                                 & 精神的健康                            & 5   &                                                     & 精神的健康                & 0.790   \\
                                 & アタッチメント                          & 9   &                                                     & 不安                   & 0.800   \\
                                 &                                  & 9   &                                                     & 回避                   & 0.740   \\
                                 & 孤独感                              & 20  &                                                     & 孤独感                  & 0.910   \\
                                 & 精神的健康                            & 5   &                                                     & 精神的健康                & 0.850   \\
        \citet{Sugiyama202192}   & 孤独感                              & 20  &                                                     & 孤独感                  &         \\
                                 & 不安・抑うつ                           & 14  &                                                     & 不安・抑うつ               &         \\
                                 & 睡眠の質                             & 9   &                                                     & 睡眠の質                 &         \\
                                 & クロノタイプ                           & 5   &                                                     & クロノタイプ               &         \\
                                 & 孤独感                              & 20  &                                                     & 孤独感                  &         \\
                                 & 不安・抑うつ                           & 14  &                                                     & 不安・抑うつ               &         \\
                                 & 睡眠の質                             & 9   &                                                     & 睡眠の質                 &         \\
                                 & クロノタイプ                           & 5   &                                                     & クロノタイプ               &         \\
                                 & 孤独感                              & 20  &                                                     & 孤独感                  &         \\
                                 & 不安・抑うつ                           & 14  &                                                     & 不安・抑うつ               &         \\
                                 & 睡眠の質                             & 9   &                                                     & 睡眠の質                 &         \\
                                 & クロノタイプ                           & 5   &                                                     & クロノタイプ               &         \\
        \citet{Kousaka202192}    & 臨時休業期間中の生活週間の変化                  & 7   & 就寝時間が不規則になった                                        & 不規則な睡眠               & 0.810   \\
                                 &                                  & 7   & 食事の量が変わった（増えた/減った）                                  & 食習慣の乱れ               & 0.760   \\
                                 &                                  & 7   & いつもより勉強する量が少なかった                                    & 学習時間の減少              & 0.810   \\
                                 &                                  & 7   & 外出する機会が減った                                          & 外出の減少                & 0.710   \\
                                 &                                  & 7   & テレビやインターネットを見ている時間が増えた                              & テレビ・ネット視聴の増加         & 0.850   \\
                                 &                                  & 7   & これまでよりも運動しなくなった                                     & 運動の減少                & 0.670   \\
                                 &                                  & 7   & ゲームやスマートフォンをいつもよりも長く利用していた                          & ゲーム・スマホ利用の増加         & 0.740   \\
        \citet{Hirai202192}      & 家庭と仕事における調整                      & 10  & 職場でお互いの事情について説明しあう                                  & 仕事調整                 & 0.822   \\
                                 &                                  & 10  & 夫婦でお互いの役割分担について話し合う                                 & 家庭調整                 & 0.806   \\
                                 & 「家族する」概念                         & 8   & 家族が自分にどうして欲しいかを考えて行動する                              & 家族する                 & 0.830   \\
                                 &                                  & 8   & 家族が自分にどうして欲しいかを考えて行動する                              & 家族する                 & 0.880   \\
                                 & 仕事へのコミットメント                      & 8   & 仕事においては人並み以上に努力をしている                                & 仕事へのコミットメント          & 0.920   \\
                                 &                                  & 8   & 仕事においては人並み以上に努力をしている                                & 仕事へのコミットメント          & 0.930   \\
                                 & 家庭，仕事，および生活・人生の満足度               & 8   & 現在の夫婦の関係                                            & 家庭                   & 0.900   \\
                                 &                                  & 8   & 現在の仕事の内容                                            & 仕事                   & 0.830   \\
                                 &                                  & 8   & 現在の自分の生活                                            & 人生                   & 0.840   \\
        \citet{Koiwa202192}      & 新型コロナウイルス感染症に対する対処行動             & 10  & 換気が悪い場所には行かないようにした                                  & 三密の回避                & 0.850   \\
                                 &                                  & 10  & 手洗い・うがいやアルコールによる手や指を消毒した                            & 日々の感染予防              & 0.630   \\
                                 &                                  & 10  & 気を紛らわすようなことをした                                      & 不安からの逃避              & 0.480   \\
                                 &                                  &     &                                                     &                      & 0.870   \\
        \citet{Shimizu202292}    & 象徴的障害者偏見尺度                       & 13  & 障害者は一生懸命努力しても，目標をたいてい達成できない                         & 個人主義                 &         \\
                                 &                                  & 13  & 日本において，障害者に対する差別はもはや問題事ではない                         & 現状の理解の無さ             &         \\
                                 & 構成世界信念                           & 12  & 不公平に苦しむ全ての人々が報われる日がいつか来るに違いない                       & 究極的構成世界信念            & 0.880   \\
                                 &                                  & 12  & どんな人であっても自分の働いた悪事の報いはいつか受けるものである                    & 内在的構成世界信念            & 0.920   \\
                                 &                                  & 12  & 世の中の大抵のことは不公平だ                                      & 不公正世界信念              & 0.830   \\
                                 & 能力の認知                            & 3   & 障害者は有能である                                           &                      & 0.850   \\
                                 & 保護の対象としての認知                      & 1   & 障害者は守られるべきである                                       &                      &         \\
                                 & 第三者公正感受性                         & 10  & 他の人より不当に苦しい生活を送っている人がいると，落ち着かない気分になる                &                      & 0.850   \\
                                 & 弱者救済規範意識                         & 5   & 虐げられている人を，まず敬うべきだ                                   &                      & 0.790   \\
                                 & 傲慢さの認知                           & 1   & 障害者はわがままである                                         &                      &         \\
        \citet{Toyama202292}     & 困難な目標への対処法略尺度                    & 25  & その目標のことは考えないようにする                                   & 目標断念                 & 0.870   \\
                                 &                                  & 25  & 新しい他の目標にやりがいを見出す                                    & 目標の内容の調整             & 0.880   \\
                                 &                                  & 25  & 最初に設定した目標よりも少しだけレベルを下げて取り組む                         & 目標の水準の調整             & 0.900   \\
                                 &                                  & 25  & その目標を達成数r為に他に良い方法はないか考える                            & 目標達成宝暦の調整            & 0.890   \\
                                 &                                  & 25  & 達成することができないとしても，その目標を追求し続ける                         & 目標継続                 & 0.880   \\
                                 &                                  & 25  & その目標のことは考えないようにする                                   & 目標断念                 & 0.800   \\
                                 &                                  & 25  & 新しい他の目標にやりがいを見出す                                    & 目標の内容の調整             & 0.810   \\
                                 &                                  & 25  & 最初に設定した目標よりも少しだけレベルを下げて取り組む                         & 目標の水準の調整             & 0.820   \\
                                 &                                  & 25  & その目標を達成数r為に他に良い方法はないか考える                            & 目標達成宝暦の調整            & 0.810   \\
                                 &                                  & 25  & 達成することができないとしても，その目標を追求し続ける                         & 目標継続                 & 0.810   \\
                                 & ホープ尺度                            & 8   & 窮地(困難な場面) から逃れる多くの方法を考えつくことができる                     & ホープ尺度                & 0.740   \\
                                 & 努力の粘り強さ                          & 6   & 私は頑張り屋だ                                             & 努力の粘り強さ              & 0.770   \\
                                 & 認知の柔軟性尺度                         & 10  & 困難な状況について対応する前に，複数の選択肢を考慮する                         & 代替                   & 0.850   \\
                                 & 適応的諦観尺度                          & 10  & 失敗してもなんとかやっていける気がする                                 &                      & 0.890   \\
                                 & 主観的well-being                    & 9   &                                                     & 抑うつ                  & 0.850   \\
                                 &                                  & 9   &                                                     & 不安                   & 0.830   \\
                                 &                                  & 9   &                                                     & 人生に対する満足             & 0.890   \\
                                 & 心理的well-being                    & 12  &                                                     & 人格的成長                & 0.840   \\
                                 &                                  & 12  &                                                     & 人生における目的             & 0.870   \\
                                 &                                  & 12  &                                                     & 自己受容                 & 0.900   \\
                                 &                                  & 12  &                                                     & 環境制御力                & 0.890   \\
                                 & 困難な目標への対処法略尺度                    & 25  & その目標のことは考えないようにする                                   & 目標断念                 & 0.820   \\
                                 &                                  & 25  & 新しい他の目標にやりがいを見出す                                    & 目標の内容の調整             & 0.830   \\
                                 &                                  & 25  & 最初に設定した目標よりも少しだけレベルを下げて取り組む                         & 目標の水準の調整             & 0.830   \\
                                 &                                  & 25  & その目標を達成する為に他に良い方法はないか考える                            & 目標達成宝暦の調整            & 0.810   \\
                                 &                                  & 25  & 達成することができないとしても，その目標を追求し続ける                         & 目標継続                 & 0.840   \\
        \citet{Nagatani202292}   & 日本語版BRIEF-T                      & 86  &                                                     & 抑制                   & 0.930   \\
                                 &                                  & 86  &                                                     & シフト                  & 0.860   \\
                                 &                                  & 86  &                                                     & 情緒コントロール             & 0.640   \\
                                 &                                  & 86  &                                                     & 開始                   & 0.900   \\
                                 &                                  & 86  &                                                     & ワーキングメモリ             & 0.890   \\
                                 &                                  & 86  &                                                     & 計画/組織化               & 0.920   \\
                                 &                                  & 86  &                                                     & 整理                   & 0.910   \\
                                 &                                  & 86  &                                                     & モニタ                  & 0.920   \\
                                 &                                  & 86  &                                                     & 抑制                   & 0.920   \\
                                 &                                  & 86  &                                                     & シフト                  & 0.880   \\
                                 &                                  & 86  &                                                     & 情緒コントロール             & 0.900   \\
                                 &                                  & 86  &                                                     & 開始                   & 0.810   \\
                                 &                                  & 86  &                                                     & ワーキングメモリ             & 0.880   \\
                                 &                                  & 86  &                                                     & 計画/組織化               & 0.880   \\
                                 &                                  & 86  &                                                     & 整理                   & 0.920   \\
                                 &                                  & 86  &                                                     & モニタ                  & 0.890   \\
                                 & 日本語版ADHD-RS                      & 18  &                                                     & 不注意尺度                &         \\
                                 &                                  & 18  &                                                     & 多動・衝動性尺度             &         \\
        \citet{Li202292}         & Promotion/Prevention Focus Scale & 16  &                                                     & 利得接近志向尺度             & 0.838   \\
                                 &                                  & 16  &                                                     & 損失回避志向尺度             & 0.857   \\ \hline
    \end{longtable}
\end{landscape}